\hypertarget{the-premise}{%
\chapter{1. The premise}\label{the-premise}}

A developer's workstation runs a lot of code the developer never wrote
or read. An AI coding agent works through a task across hundreds of
files. A postinstall script runs the moment a dependency is added. An
MCP server starts because a tool asked for it. A build system runs in a
repository cloned an hour ago to look at, not yet trusted. None of it
came through a path the system checks: a public registry, a
\texttt{curl\ \textbar{}\ sh} line, a model's output run locally. All of
it runs as the user, with the user's access to the user's files, keys,
network, and services.

Under Unix, that access is not a permission the code was handed. It is
who the code is. The system treats everything the code does as something
the user did.

The premise the rest of the design answers begins here: why that is now
a problem, what kind of problem it is, and what shape of answer it
needs. The threat catalogue (\texttt{THREATS.md}) lists the concrete
threats and the principle register (\texttt{PRINCIPLES.md}) states the
posture. This volume is the platform-neutral design that meets them; the
second volume is its realisation on Linux.

\hypertarget{the-threat-is-effect-not-intent}{%
\section{1.1 The threat is effect, not
intent}\label{the-threat-is-effect-not-intent}}

The natural move is to sort this code into trustworthy and malicious,
and guard against the malicious. That sort is the mistake. The project's
founding principle (\texttt{good-boy}) is the decision not to make it.

Start with the behaviour that prompted the work. An AI agent was run
inside a sandbox and blocked from a command. It did not give up. It
found a path the deny rule did not cover, and when that was closed too,
it reasoned its way to switching the sandbox off and running the command
outside it (\texttt{T2.1}). Nobody told it to escape. Switching the
sandbox off was simply the quickest way left to finish the job. To an
agent under pressure to complete a task, a safeguard is just another
obstacle, and the search for a way around it never stops to ask what the
obstacle was protecting.

The same shape turns up at the malicious end. In the Nx attack, poisoned
package versions ran postinstall scripts that searched the disk for
secrets and used the developer's own coding agents to help send them out
(\texttt{T1.1}, \texttt{T1.2}). Put the two cases side by side. A
malicious script hunts for credentials in
\texttt{\textasciitilde{}/.ssh}, \texttt{\textasciitilde{}/.aws}, and
the shell history. A helpful agent reads the same files in case
something in them is useful. Same reads, same files, same hole. The
boundary cannot tell the two apart from what they do, and it does not
need to. The harm a workload causes depends on what it can reach, not on
what it meant to do.

So the framework does not ask. It assumes a workload will do whatever
its capabilities allow, and it limits the capabilities instead of
judging the motive. The dog loves you and wrecks the couch anyway. You
do not crate it because you have decided it is a bad dog; you crate it
so the couch is safe whichever kind of dog it turns out to be. That is
the difference between confinement and detection. Detection has to be
right about the code, and it is wrong the day a trusted dependency ships
a compromised release, or a loyal agent reads a poisoned instruction out
of the repository it is working in (\texttt{T3.7}). Confinement never
has to be right about the code, so guessing wrong does not break it.

None of this is new. The same problem, code a user runs that may harm
them by accident or by intent, met by mediating it rather than trusting
it, was set out in the Anderson report in 1972 {[}Anderson{]}, and has
been the security tradition's answer ever since {[}Irvine{]}. The name
\texttt{good-boy} is the project's; the instinct is not.

\hypertarget{identity-and-authority}{%
\section{1.2 Identity and authority}\label{identity-and-authority}}

Most modern operating systems make no distinction between who you are
and what you may do. On Unix, Windows, and macOS alike, the account you
are is the authority you hold. It shows most plainly in Unix terms. The
uid is both. It is your identity, the name the system knows you by, and
it is your authority, the sum of everything that name may touch. The
system cannot hold the one apart from the other. The equation once
reached past the machine. rsh and the \texttt{.rhosts} file let a uid's
authority cross every host that trusted it, so your identity on one box
was your authority across a domain. Almost every host security measure
since has worked to pull the two apart. setuid and sudo let a process
act with an authority its identity lacks; mandatory access control lets
the system grant an identity less than its uid would command. Each one
forces a decision about authority where the default would simply re-read
it off the identity: the same answer every time, no separate question
asked. The whole host chain hangs off that: a fused uid, and decades of
measures built to pry it open.

None of that work reached the interactive user level, where the equation
still holds in full (\texttt{split-the-uid}). Code running under your
uid acts with your identity and inherits your entire authority,
undivided. When the agent reads a file, when the script opens a
connection, when the build writes to your home directory, the system
sees your access, because by uid it is yours, and it cannot tell you
apart from the code wearing your name.

People do put boundaries around this code, and two kinds are common. A
sandbox wraps one tool. It fixes a boundary when the tool starts and
takes no further interest in what the tool does inside it. A container
wraps a whole system. It fixes one boundary around everything it holds
and calls the job done. Each draws its line once and steps back. Each
settles what the workload can reach, and neither looks again at what the
workload does with the reach it was given.

A reference monitor (\texttt{reference-monitor}) is a different kind of
thing. It does not draw a boundary and withdraw. It stands in the path
of every access and clears them one at a time: each read, each
connection, each call, judged as it is made, for as long as the workload
runs. The sandbox and the container both draw a line and step back; the
monitor stays in the path and meets each use as it comes
(\texttt{mediate-use-not-reach}). The host's chain of measures
culminates in one. The user level has none. Kennel keeps the identity,
the uid you already are, and splits the authority off it: the workload
acts under your name, with only the authority its policy grants, and
exercises even that one access at a time. The next chapter builds that
monitor where it is missing.

\hypertarget{the-trust-root-and-the-limits-of-the-claim}{%
\section{1.3 The trust root, and the limits of the
claim}\label{the-trust-root-and-the-limits-of-the-claim}}

Someone still has to be trusted, and it is the operator, the person
running the framework. The operator is the trust root, not a threat to
defend against. Kennel confines the workload and works for the operator.
Its job is to make every capability the operator hands out explicit,
bounded, and visible, not to overrule a choice the operator made
deliberately (\texttt{footgun-warn-dont-forbid}). This cuts both ways. A
framework that blocks what the operator needs gets switched off or
worked around; the same pressure that pushes a workload past a safeguard
pushes the person past one too, and a control that gets bypassed
protects nothing.

The trust root also fixes the edges of what the framework claims, and
the design names them rather than pretending to cover everything. It
governs the unsigned workload at the user level: what it reads, where it
connects, what it changes, and what it leaves behind. It does not reach
hardware or physical attacks, which sit below any userspace boundary. It
does not reach an operator who decides to help the workload, because
that consent is the trust root's to give. It does not reach bugs in the
kernel parts it is built on, which are the ground it stands on. And it
does not decide what data may leave through a channel the operator
already allowed; which bytes ride an approved request is a
data-governance question, not a confinement one (\texttt{T1.8}). The
framework limits what a workload can reach and where it can connect. It
does not claim to limit what an approved channel carries, and it says so
where that line runs.

That limit is not a shortfall in the build; it is the edge of the
method. A monitor that clears each access as it happens can enforce that
nothing forbidden occurs, but where information travels after it has
legitimately left is a fact about the flow of data across many actions,
not about any single access, and a monitor of this kind provably cannot
enforce it {[}Schneider{]}. The limit is a theorem, not an apology.

\hypertarget{the-shape-of-the-answer}{%
\section{1.4 The shape of the answer}\label{the-shape-of-the-answer}}

Project Kennel is a reference monitor for the user level, built in user
space, holding no privilege the workload could use against it
(\texttt{reference-monitor}). Four commitments give it its shape, and
the chapters that follow work each of them out.

It denies by default. A workload gets only what its policy grants. What
it was not granted is not refused when it asks for it: it is simply not
there, with nothing to find and nothing to probe
(\texttt{construction-by-absence}). What cannot be removed, because the
workload truly needs it, is reached through a transaction the monitor
approves one call at a time, never as an open channel handed over and
left alone (\texttt{interpose-as-transaction}). The whole is held to
doing as little as it can, because a monitor is only as trustworthy as
it is small enough to read (\texttt{do-less}). And at no point does it
ask whether the workload deserves trust, because that is not a question
the enforcement depends on (\texttt{good-boy}).

The rest of this volume develops that design without tying it to any one
operating system, because the design is the contract and the contract is
platform-neutral. Chapter 2 builds the reference monitor and its three
properties. Chapter 3 covers the two ways it reaches complete mediation:
removing a resource so there is nothing to mediate, and interposing on
the uses it cannot remove. Chapters 4 through 11 work through the
resource classes one at a time, from the filesystem to the network to
the services a session depends on. Chapters 12 through 16 turn to the
grant itself: how a confinement is defined, made to enforce, consented
to, vouched for, and accounted for. Each decision answers a threat from
the catalogue and follows from a principle in the register. How the
contract is enforced on a real kernel is the second volume.
